\section{Node.js}\label{sec:node_js}

Node.js adalah lingkungan eksekusi (\textit{runtime environment}) untuk JavaScript yang bersifat \textit{open-source} dan lintas platform, dibangun di atas mesin JavaScript V8 milik Chrome. Node.js menggunakan model I/O yang bersifat \textit{non-blocking} dan \textit{event-driven} yang menjadikannya ringan dan efisien, terutama untuk aplikasi yang berjalan secara \textit{real-time} pada perangkat terdistribusi (Tilkov \& Vinoski, 2010).

Dalam konteks penelitian ini, Node.js dipilih sebagai platform utama pengembangan sistem {data capture} dan integrasi model. Pemilihan ini didasarkan pada kebutuhan kompatibilitas dengan sistem \texttt{srusun.id} yang digunakan oleh pengelola Apartemen The Jarrdin, serta untuk memastikan integrasi yang mulus dengan modul-modul lain yang dikembangkan dalam ekosistem yang sama. Node.js memungkinkan penanganan permintaan data penyewa secara asinkronus, sehingga proses input dan validasi data penyewa dapat dilakukan dengan cepat tanpa membebani kinerja server secara keseluruhan.

Untuk mendukung fungsionalitas sistem, Node.js didukung oleh ekosistem pustaka yang luas melalui NPM (\textit{Node Package Manager}). Beberapa pustaka utama yang digunakan dalam penelitian ini meliputi \texttt{express} untuk pengelolaan server dan rute API, \texttt{mysql2/promise} untuk koneksi asinkronus ke basis data melalui mekanisme \textit{connection pool}, \texttt{jsonwebtoken} untuk keamanan autentikasi lintas platform, \texttt{multer} untuk menangani unggahan berkas (\textit{file upload}), \texttt{json2csv} untuk keperluan ekspor data laporan, serta \texttt{cors} untuk menangani akses lintas domain. Berikut langkah-langkah implementasi sistem \textit{backend} menggunakan Node.js.

\begin{enumerate}
    \item \textbf{Persiapan lingkungan dan instalasi pustaka}
    \begin{enumerate}
        \item Menginisialisasi proyek Node.js dan menginstal pustaka yang dibutuhkan untuk server dan basis data.
        \item Menginstal dependensi utama melalui terminal.
        \begin{lstlisting}[language=bash]
npm init -y
npm install express mysql2 dotenv path body-parser cors multer json2csv jsonwebtoken
        \end{lstlisting}
        \item Menyiapkan variabel lingkungan (\textit{environment variables}) pada berkas \texttt{.env} untuk menyimpan konfigurasi basis data jarak jauh (\textit{remote}) dan kunci rahasia secara aman.
    \end{enumerate}
        
    % \newpage
    \item \textbf{Konfigurasi server dan koneksi basis data}
    \begin{enumerate}
        \item Membuat konfigurasi server menggunakan \textit{framework} Express.js dan mengatur \textit{middleware} untuk keamanan dan format pertukaran data.
        \begin{lstlisting}[language=JavaScript, caption={Inisialisasi Server dan \textit{Middleware}}]
const express = require('express');
const cors = require('cors');
const bodyParser = require('body-parser');
const jwt = require('jsonwebtoken');

const app = express();

const corsOptions = {
    origin: (origin, callback) => {
        // Validasi domain yang diizinkan
        callback(null, true); 
    },
    credentials: true,
    methods: ['GET', 'POST', 'PUT', 'PATCH', 'DELETE', 'OPTIONS'],
    allowedHeaders: ['Content-Type', 'Authorization']
};

app.use(cors(corsOptions));
app.use(bodyParser.json());
app.use(bodyParser.urlencoded({ extended: true }));
        \end{lstlisting}
             
        % \newpage
        \item Membangun koneksi terpusat ke basis data MySQL (SRusun) menggunakan \textit{Connection Pool} untuk mengelola lalu lintas kueri secara efisien.
        \begin{lstlisting}[language=JavaScript, caption={Konfigurasi \textit{Connection Pool} Basis Data}]
const mysql = require('mysql2/promise');
require('dotenv').config();

const pool = mysql.createPool({
    host: process.env.SRUSUN_MYSQL_HOST,
    user: process.env.SRUSUN_MYSQL_USER,
    password: process.env.SRUSUN_MYSQL_PASSWORD,
    database: process.env.SRUSUN_MYSQL_DATABASE,
    waitForConnections: true,
    connectionLimit: 10,
    queueLimit: 0
});
        \end{lstlisting}
    \end{enumerate}
        
    \item \textbf{Implementasi API {data capture}} \\
    Membuat \textit{endpoint} API (\texttt{POST /api/rentals}) untuk menerima input data penyewa dari antarmuka pengguna (\textit{frontend}). Implementasi ini menggunakan mekanisme \texttt{\textit{database transaction}} untuk menjamin integritas data (ACID), di mana penyimpanan data profil penyewa dan pencatatan riwayat aktivitas (\textit{audit trail}) dilakukan dalam satu kesatuan proses yang saling bergantung.

    \begin{lstlisting}[language=JavaScript, caption={Implementasi Transaksi API Data Capture}]
app.post('/api/rentals', checkAuth, async (req, res) => {
    const conn = await pool.getConnection();
    try {
        const {
            nama, nik, status_pasutri, status_kewarganegaraan, 
            jenis_sewa, unit_number, metode_pembayaran, metode_lain, 
            tanggal_checkin, waktu_checkin, tanggal_checkout
        } = req.body;

        // Normalisasi NIK dan penentuan agen
        const normalizedNik = String(nik || '').replace(/\D/g, '');
        let agent_email = req.user.email;
        if (isAdmin(req.user) && req.body.agent_email) {
            agent_email = req.body.agent_email;
        }

        const lama_menginap = calculateDays(tanggal_checkin, tanggal_checkout);
        const db_status_pasutri = status_pasutri === 'Ya' ? 'Menikah' : 'Belum Menikah';

        // Memulai transaksi basis data
        await conn.beginTransaction();

        // Menyimpan data profil ke tabel rentals
        const sql = `INSERT INTO rentals (
            nama, nik, status_pasutri, status_kewarganegaraan, 
            jenis_sewa, unit_number, metode_pembayaran, metode_lain, 
            tanggal_checkin, waktu_checkin, tanggal_checkout, lama_menginap, user_email
        ) VALUES (?, ?, ?, ?, ?, ?, ?, ?, ?, ?, ?, ?, ?)`;

        const params = [
            nama, normalizedNik, db_status_pasutri, status_kewarganegaraan,
            jenis_sewa, unit_number, metode_pembayaran, 
            metode_pembayaran === 'Others' ? metode_lain : null,
            tanggal_checkin, waktu_checkin, tanggal_checkout, lama_menginap, agent_email
        ];

        const [result] = await conn.query(sql, params);
        const newRentalId = result.insertId;

        // Menyimpan rekam jejak pembuatan awal ke tabel rental_logs
        const logQuery = `INSERT INTO rental_logs (action, field_changed, new_value, email, rental_id) VALUES ?`;
        const logEntries = [['create', 'all', JSON.stringify(req.body), req.user.email, newRentalId]];
        await conn.query(logQuery, [logEntries]);

        // Menyelesaikan transaksi jika seluruh kueri berhasil
        await conn.commit();
        res.json({ success: true, message: "Data penyewa berhasil disimpan." });

    } catch (error) {
        // Membatalkan seluruh perubahan jika terjadi galat pada salah satu tahapan
        if (conn) await conn.rollback();
        console.error("Error saving rental:", error);
        res.status(500).json({ success: false, message: 'Gagal menyimpan data penyewa.' });
    } finally {
        if (conn) conn.release();
    }
});
    \end{lstlisting}
\end{enumerate}

Data yang telah tersimpan secara terstruktur dalam basis data SRusun kemudian diakses dan diolah lebih lanjut menggunakan bahasa pemrograman Python untuk tahapan klasifikasi.
